\section{Student Learning and Responsible Use}\label{sec:student}

AI can support personalized, accessible, and engaging learning, while also posing risks of over-reliance, inaccuracy, and bias \citep{kasneci2023}. Studies of student perspectives in higher education report both perceived benefits and concerns, which supports a guided rather than prohibitive approach \citep{chanhu2023}.

\subsection{Student expectations}
Students are expected to:
\begin{enumerate}
  \item read each course's AI-use declaration and follow it; ask the instructor when unclear;
  \item use AI to deepen understanding, not to bypass the learning the course intends;
  \item verify AI outputs against reliable sources; they remain responsible for the accuracy of submitted work, including fabricated or incorrect citations;
  \item disclose AI use when the course requires it, using a short statement (Appendix~\ref{app:declarations});
  \item protect their own and others' personal data when using AI tools; and
  \item keep drafts, notes, or AI-chat records that show how work developed, where the course asks for process evidence.
\end{enumerate}

\subsection{Learning-enhancing uses}
Where a course permits, examples of encouraged student uses include: explaining difficult concepts, generating practice questions, getting feedback on drafts, language support, debugging and code explanation, brainstorming, summarizing readings (followed by reading the source), and organizing study plans.

\subsection{Guard against over-reliance}
Faculty are encouraged to design activities in which students critique, compare, and improve AI outputs, and to build reflection on how AI use affected learning. Students are encouraged to attempt tasks first before consulting AI when building foundational skills.

\subsection{Student voice}
Students are represented in the Steering Committee and in the feedback loop. Annual surveys and focus groups inform revisions.
